\section{Reproducibility}

\subsection{Environment}

Base dependencies are listed in \texttt{requirements.txt}: \texttt{pandas}, \texttt{numpy}, \texttt{pillow}, \texttt{tqdm}, \texttt{requests}, \texttt{hydra-core}, \texttt{torch}, and \texttt{torchvision}. Additional runtime packages required by advanced features are:
\begin{itemize}
  \item \texttt{open\_clip\_torch} for the OpenCLIP backend.
  \item \texttt{ultralyticsplus} for bundle clothing detection.
\end{itemize}

\subsection{Determinism Controls}

The code explicitly sets random seeds (\texttt{random}, \texttt{numpy}, \texttt{torch}). Offline augmentation uses stable hash-based seeds per asset and augmentation index. This enables repeatable preprocessing and consistent training behavior under fixed configuration.

\subsection{Hydra Configuration}

Configuration is split into:
\begin{itemize}
  \item \texttt{config/config.yaml}: training and inference parameters.
  \item \texttt{config/files/file.yaml}: dataset/image/cache paths.
\end{itemize}

Hydra output directories store run-specific artifacts, making experiments traceable without manual path bookkeeping.

\subsection{Submission Contract}

The final CSV must contain:
\begin{itemize}
  \item \texttt{bundle\_asset\_id}
  \item \texttt{product\_asset\_id}
\end{itemize}
with up to 15 rows per bundle, sorted by model confidence.
